% =========================================================
% ESTILOS Y COMANDOS PROPIOS DEL INFORME
% Complementa a pucv_inf_2024.sty sin modificarlo.
% =========================================================

% ---------------------------------------------------------
% Rótulos en español
% Se reaplican después de babel para conservar los títulos de
% índices definidos por pucv_inf_2024.sty.
% ---------------------------------------------------------

\AtBeginDocument{%
    \renewcommand{\figurename}{Figura}%
    \renewcommand{\tablename}{Tabla}%
    \renewcommand{\nomname}{\uppercase{Lista de Símbolos}}%
    \renewcommand{\contentsname}{\uppercase{Índice General}}%
    \renewcommand{\listfigurename}{\uppercase{Lista de Figuras}}%
    \renewcommand{\listtablename}{\uppercase{Lista de Tablas}}%
    \renewcommand{\chaptername}{Capítulo}%
    \renewcommand{\appendixname}{Anexo}%
    \renewcommand{\partname}{Parte}%
}

% ---------------------------------------------------------
% Bibliografía en español
% Fuerza los rótulos APA en español («Recuperado el ... de»,
% «2.ª ed.», «y» en lugar de «and») sin depender de babel.
% ---------------------------------------------------------

\ExecuteBibliographyOptions{language = spanish}
\DeclareLanguageMapping{spanish}{spanish-apa}

% ---------------------------------------------------------
% Colores auxiliares
% ---------------------------------------------------------

\definecolor{notacolor}{HTML}{4472c4}   % azul institucional (igual a titlecolor)
\definecolor{notafondo}{HTML}{EEF2FA}
\definecolor{pendientecolor}{HTML}{B85C00}
\definecolor{pendientefondo}{HTML}{FDF3E7}
\definecolor{filacabecera}{HTML}{DCE6F5}

% ---------------------------------------------------------
% Cajas de nota y de tareas pendientes
% (equivalen a las citas en bloque y a los bloques
%  "POR COMPLETAR" del documento original en Markdown)
% ---------------------------------------------------------

\newtcolorbox{nota}[1][]{%
    enhanced,
    breakable,
    colback = notafondo,
    colframe = notacolor,
    boxrule = 0pt,
    leftrule = 3pt,
    arc = 0pt,
    outer arc = 0pt,
    left = 8pt, right = 8pt, top = 6pt, bottom = 6pt,
    fontupper = \small,
    parbox = false,
    #1
}

\newtcolorbox{pendiente}[1][Por completar]{%
    enhanced,
    breakable,
    colback = pendientefondo,
    colframe = pendientecolor,
    boxrule = 0.6pt,
    arc = 2pt,
    left = 8pt, right = 8pt, top = 6pt, bottom = 6pt,
    fontupper = \small,
    parbox = false,
    title = {\footnotesize\bfseries\MakeUppercase{#1}},
    coltitle = white,
    colbacktitle = pendientecolor
}

% Marca breve, para usar dentro de una línea de texto
\newcommand{\porcompletarinline}[1]{%
    \textcolor{pendientecolor}{\textbf{[por completar: #1]}}%
}

% ---------------------------------------------------------
% Estilo común de las tablas largas
% ---------------------------------------------------------

\newcommand{\estilotabla}{%
    \footnotesize
    \setlength{\parskip}{0pt}%
    \setlength{\parindent}{0pt}%
    \setlength{\tabcolsep}{4pt}%
    \renewcommand{\arraystretch}{1.25}%
}

\newcommand{\estilotablachica}{%
    \estilotabla
    \scriptsize
}

% Separación antes y después de las longtable
\setlength{\LTpre}{8pt}
\setlength{\LTpost}{8pt}

% Columnas de párrafo centradas verticalmente arriba
\newcolumntype{L}[1]{>{\raggedright\arraybackslash}p{#1}}
\newcolumntype{C}[1]{>{\centering\arraybackslash}p{#1}}

% ---------------------------------------------------------
% Identificadores de código, clases y estados
% ---------------------------------------------------------

\newcommand{\clase}[1]{\texttt{#1}}
\newcommand{\estado}[1]{\texttt{#1}}
\newcommand{\este}[1]{\textit{#1}}

% Estereotipos UML
\newcommand{\ster}[1]{\texttt{\guillemotleft #1\guillemotright}}

% ---------------------------------------------------------
% Figuras de diagramas: comando abreviado
% #1 ancho relativo · #2 archivo · #3 título breve (Lista de Figuras)
% #4 leyenda completa · #5 etiqueta
% ---------------------------------------------------------

\newcommand{\figuradiagrama}[5]{%
    \begin{figure}[H]
        \centering
        \IfFileExists{Figuras/diagramas/#2}{%
            \includegraphics[width=#1\textwidth]{#2}%
        }{%
            \fbox{\parbox[c][7cm][c]{0.85\textwidth}{\centering
                \textbf{[Diagrama pendiente]}\\[4pt]
                \texttt{Figuras/diagramas/#2}}}%
        }
        \caption[#3]{#4}
        \label{#5}
    \end{figure}%
}

% ---------------------------------------------------------
% Tablas de casos de uso narrativos (formato de los ejemplos)
% ---------------------------------------------------------

\newcommand{\campocu}[2]{\hline \textbf{#1} & #2 \\}
\newcommand{\encabezadocu}[1]{\hline \rowcolor{filacabecera} \multicolumn{2}{|c|}{\textbf{#1}} \\}
\newcommand{\subencabezadocu}[2]{\hline \rowcolor{filacabecera} \textbf{#1} & \textbf{#2} \\}

% ---------------------------------------------------------
% Ajustes de listas y espaciado dentro de entornos
% ---------------------------------------------------------

\setlist[itemize]{leftmargin=1.2cm, itemsep=2pt, topsep=4pt}
\setlist[enumerate]{leftmargin=1.2cm, itemsep=2pt, topsep=4pt}

% Evita que hyperref rompa las referencias de longtable
\hypersetup{
    bookmarksnumbered = true,
    pdftitle = {Informe 3 - Análisis y diseño estructurado del sistema BeeTracer},
    pdfsubject = {ICI-3242 Modelamiento de Software - PUCV}
}

% Nombre de botón o menú entre comillas angulares
\newcommand{\bt}[1]{\guillemotleft{}#1\guillemotright{}}

% Numeración corrida de figuras y tablas (Figura 1, 2, 3... como en los ejemplos)
\usepackage{chngcntr}
\counterwithout{figure}{chapter}
\counterwithout{table}{chapter}

% ---------------------------------------------------------
% Tolerancia de justificación: evita desbordes de línea con
% nombres de clases largos en monoespaciado.
% ---------------------------------------------------------

\setlength{\emergencystretch}{3em}

% ---------------------------------------------------------
% Páginas de "Parte" (Parte 1: Análisis · Parte 2: Diseño)
% ---------------------------------------------------------

\renewcommand{\thepart}{\arabic{part}}
\titleformat{\part}[display]
    {\centering\normalfont\bfseries\Huge\color{titlecolor}}
    {\MakeUppercase{\partname}~\thepart}{1em}{\MakeUppercase}

% ---------------------------------------------------------
% Estilos TikZ para los diagramas del análisis y diseño
% estructurado (notación de Yourdon/DeMarco, Chen y
% diagramas de estructura de Yourdon y Constantine)
% ---------------------------------------------------------

\usetikzlibrary{arrows.meta, positioning, calc, fit, backgrounds,
    shapes.geometric, shapes.multipart, decorations.markings}

\definecolor{dfdproceso}{HTML}{EAF1FB}
\definecolor{dfdterminador}{HTML}{E6E6E6}
\definecolor{dfdalmacen}{HTML}{FFF8E6}
\definecolor{dfdborde}{HTML}{2F4F7F}

\tikzset{
    % --- Diagramas de flujo de datos ---
    proceso/.style = {circle, draw = dfdborde, line width = 0.7pt,
        fill = dfdproceso, align = center, font = \scriptsize,
        minimum size = #1, inner sep = 1pt},
    proceso/.default = 2.3cm,
    terminador/.style = {rectangle, draw = black, line width = 0.7pt,
        fill = dfdterminador, align = center, font = \scriptsize\bfseries,
        minimum width = 2.5cm, minimum height = 0.95cm, inner sep = 3pt},
    almacen/.style = {rectangle, fill = dfdalmacen, align = center,
        font = \scriptsize, minimum width = #1, minimum height = 0.6cm,
        inner sep = 2pt, outer sep = 0pt,
        append after command = {
            (\tikzlastnode.north west) edge[draw, -, line width = 0.8pt]
            (\tikzlastnode.north east)
            (\tikzlastnode.south west) edge[draw, -, line width = 0.8pt]
            (\tikzlastnode.south east)}},
    almacen/.default = 2.6cm,
    flujo/.style = {-{Stealth[length = 2mm]}, line width = 0.5pt},
    etq/.style = {font = \tiny, fill = white, inner sep = 1pt,
        align = center, text = black},
    % --- Diagrama entidad-relación (Chen) ---
    entidad/.style = {rectangle split, rectangle split parts = 2,
        draw = black, line width = 0.7pt,
        font = \scriptsize, inner sep = 2.5pt, minimum width = 3.0cm, align = left,
        rectangle split part align = {center, left},
        rectangle split part fill = {dfdborde!22, white}},
    relacion/.style = {diamond, aspect = 2.4, draw = black,
        line width = 0.6pt, fill = dfdalmacen, font = \scriptsize,
        align = center, inner sep = 0.5pt},
    card/.style = {font = \scriptsize\bfseries, inner sep = 1pt, fill = white},
    % --- Modelo relacional ---
    tabla/.style = {rectangle split, rectangle split parts = 2,
        draw = black, line width = 0.6pt, font = \tiny, align = left,
        inner sep = 2pt, rectangle split part fill = {dfdborde!30, white},
        rectangle split part align = {center, left}},
    fk/.style = {-{Stealth[length = 1.6mm]}, line width = 0.45pt,
        draw = dfdborde},
    % --- Diagramas de estructura ---
    modulo/.style = {rectangle, draw = black, line width = 0.7pt,
        fill = dfdproceso, align = center, font = \scriptsize,
        text width = #1, minimum height = 0.9cm, inner sep = 2pt,
        execute at begin node = {\hyphenpenalty = 10000 \exhyphenpenalty = 10000}},
    modulo/.default = 2.1cm,
    predefinido/.style = {modulo = #1,
        append after command = {
            ($(\tikzlastnode.north west) + (0.12, 0)$) edge[draw, -, line width = 0.5pt]
            ($(\tikzlastnode.south west) + (0.12, 0)$)
            ($(\tikzlastnode.north east) + (-0.12, 0)$) edge[draw, -, line width = 0.5pt]
            ($(\tikzlastnode.south east) + (-0.12, 0)$)}},
    predefinido/.default = 2.1cm,
    llamada/.style = {-{Stealth[length = 1.8mm]}, line width = 0.5pt},
    dato/.style = {{Circle[open, length = 1.6mm]}-{Stealth[length = 1.6mm]},
        line width = 0.4pt},
    control/.style = {{Circle[length = 1.6mm]}-{Stealth[length = 1.6mm]},
        line width = 0.4pt},
    acop/.style = {font = \tiny, inner sep = 0.8pt, align = center},
}

% Rótulo de proceso: número en negrita y nombre debajo
\newcommand{\proc}[2]{\textbf{#1}\\[-1pt]#2}

% Inclusión de un diagrama TikZ ajustado al ancho disponible
% #1 archivo en Figuras/tikz · #2 título breve · #3 leyenda · #4 etiqueta
\newcommand{\diagramatikz}[4]{%
    \begin{figure}[H]
        \centering
        \begin{adjustbox}{max width=\textwidth, max totalheight=0.80\textheight}
            \input{Figuras/tikz/#1}
        \end{adjustbox}
        \caption[#2]{#3}
        \label{#4}
    \end{figure}%
}

% Par de flujos opuestos entre un terminador (o almacén) y un proceso.
% #1 punto origen (externo) · #2 punto destino (proceso)
% #3 rótulo del flujo #1->#2 · #4 rótulo del flujo #2->#1
% #5 posición del rótulo #3 · #6 posición del rótulo #4 (0 a 1)
\newcommand{\parflujo}[6]{%
    \draw[flujo] ($(#1)!0.17cm!90:(#2)$) --
        node[etq, auto, pos = #5] {#3} ($(#2)!0.17cm!-90:(#1)$);
    \draw[flujo] ($(#2)!0.17cm!90:(#1)$) --
        node[etq, auto, pos = #6] {#4} ($(#1)!0.17cm!-90:(#2)$);
}

% Flujo de lectura y escritura sobre un almacén (doble punta)
\tikzset{flujorw/.style = {{Stealth[length = 2mm]}-{Stealth[length = 2mm]},
    line width = 0.5pt}}

% Diagrama TikZ en página apaisada
% #1 archivo · #2 título breve · #3 leyenda · #4 etiqueta
\newcommand{\diagramaapaisado}[4]{%
    \begin{landscape}
    \begin{figure}[p]
        \centering
        \begin{adjustbox}{max width=\linewidth, max totalheight=0.86\textheight}
            \input{Figuras/tikz/#1}
        \end{adjustbox}
        \caption[#2]{#3}
        \label{#4}
    \end{figure}
    \end{landscape}%
}

% Relación del DER entre dos entidades con su rombo y cardinalidades
% #1 entidad A · #2 entidad B · #3 nombre del rombo · #4 rótulo
% #5 cardinalidad junto a A · #6 cardinalidad junto a B · #7 posición del rombo
\newcommand{\relER}[7]{%
    \node[relacion] (#3) at ($(#1)!#7!(#2)$) {#4};
    \draw[line width = 0.5pt] (#1) -- node[card, pos = 0.4] {#5} (#3);
    \draw[line width = 0.5pt] (#3) -- node[card, pos = 0.6] {#6} (#2);
}

% ---------------------------------------------------------
% Mini especificaciones (español estructurado)
% ---------------------------------------------------------
\newtcolorbox{minispec}[2]{%
    enhanced, breakable, colback = white, colframe = dfdborde,
    boxrule = 0.6pt, arc = 1pt, left = 6pt, right = 6pt, top = 3pt, bottom = 3pt,
    title = {\textbf{Proceso #1:} #2}, fonttitle = \small, coltitle = white,
    colbacktitle = dfdborde, fontupper = \small, parbox = false,
    before skip = 6pt, after skip = 6pt}
\newcommand{\msdato}[2]{\noindent\textbf{#1:} #2\par}
\newcommand{\kw}[1]{\textsf{\textbf{#1}}}
\newlist{logica}{enumerate}{3}
\setlist[logica]{label = \arabic*., nosep, leftmargin = 1.6em, font = \footnotesize}
\setlist[logica, 2]{label = \alph*., leftmargin = 1.4em}
\setlist[logica, 3]{label = \roman*., leftmargin = 1.4em}
\newcommand{\fl}[1]{\textit{#1}}

% ---------------------------------------------------------
% Diagramas de estructura: llamadas y acoplamientos
% #1 tipo de acoplamiento (dato | control) · #2 módulo superior
% #3 módulo subordinado · #4 rótulo
% ---------------------------------------------------------
\newcommand{\conecta}[2]{\draw[llamada] (#1) -- (#2);}
\newcommand{\baja}[4][dato]{%
    \draw[#1] ($(#2)!0.30!(#3)!0.24cm!-90:(#3)$) --
        node[acop, auto, swap] {#4} ($(#2)!0.58!(#3)!0.24cm!-90:(#3)$);}
\newcommand{\sube}[4][dato]{%
    \draw[#1] ($(#2)!0.62!(#3)!0.24cm!90:(#3)$) --
        node[acop, auto, swap] {#4} ($(#2)!0.34!(#3)!0.24cm!90:(#3)$);}
% Rombo de selección (centro de transacción) en la base de un módulo
\tikzset{seleccion/.style = {diamond, draw, fill = white, inner sep = 0pt,
    minimum size = 0.32cm}}

% Columnas del modelo relacional
\newcommand{\pk}[1]{\uline{#1}}
\newcommand{\fk}[1]{\textit{#1} (FK)}

% Especificación detallada de un módulo del diseño
\newtcolorbox{modspec}[1]{%
    enhanced, breakable, colback = white, colframe = dfdborde!70!black,
    boxrule = 0.6pt, arc = 1pt, left = 6pt, right = 6pt, top = 3pt, bottom = 3pt,
    title = {\textbf{Módulo:} #1}, fonttitle = \small, coltitle = white,
    colbacktitle = dfdborde!70!black, fontupper = \small, parbox = false,
    before skip = 6pt, after skip = 6pt}

% Tabla de módulos de un subsistema
% #1 título de la tabla · #2 etiqueta · #3 filas
\newcommand{\tablamodulos}[3]{%
    {\estilotabla
    \begin{longtable}{|L{3.1cm}|L{2.4cm}|L{2.9cm}|L{2.6cm}|L{3.6cm}|}
    \caption{#1}\label{#2}\\
    \hline
    \rowcolor{filacabecera}
    \textbf{Módulo} & \textbf{Llamado por} & \textbf{Entradas} & \textbf{Salidas} &
    \textbf{Función} \\ \hline
    \endfirsthead
    \hline
    \rowcolor{filacabecera}
    \textbf{Módulo} & \textbf{Llamado por} & \textbf{Entradas} & \textbf{Salidas} &
    \textbf{Función} \\ \hline
    \endhead
    #3
    \end{longtable}}%
}
